\begingroup
\fancyhead[L]{\small POSTHUMOUS PAPERS AND CORRESPONDENCE}
\clearpage
\SourceLeaf{571}
\pdfbookmark[1]{Posthumous Papers VII. Tables for the Perturbations of Pallas by Jupiter}{pallas-nachlass-VII}
\begin{center}
\large POSTHUMOUS PAPERS AND CORRESPONDENCE.
\par\bigskip
\Large [VII. TABLES FOR THE PERTURBATIONS\par OF PALLAS BY JUPITER.]
\par\medskip
\normalsize [AUGUST 1816---DECEMBER 1817.]
\par\bigskip
[1.]
\par\medskip
\textsc{Gauss} to \textsc{Encke}, Göttingen, 5 August 1816.
\end{center}
\pdfbookmark[2]{[1] Gauss to Encke, 5 August 1816}{pallas-VII-letter1}
\par\noindent\hspace*{\parindent}.\ .\ .\ .\ .\quad At the same time I have another request. The whole calculation
of all the perturbations of Pallas for a single opposition always takes
a month because of the time spent sending it back and forth. I myself
(admittedly not continuously, but one seldom feels like doing it either)
spent nearly 14 days this time on the single calculation. I have thought
the matter over, considering how tables might shorten the work, and found
that with the aid of a certain table the work might indeed be finished in 3 hours.
But calculating the table will require an immense amount of work.
Above all, an estimate is needed of how
\textls[100]{much} work it would take, so that one can judge whether 6 or 8 calculators
could perhaps carry it out, even if it took years. To that end
I am sending you the data for a trial calculation. I am less interested
in the results (at least for now) than in knowing exactly how
much time you will have spent on it. This calculation would amount to
about \(\tfrac{1}{100}\) of the \textls[100]{whole} work.
\par\noindent\hspace*{\parindent}I should like the values
\SourceLeaf{572}
\[
\begin{aligned}
\text{1) for }P={}&+359\overset{\prime\prime}{,}834\\
&+54{,}481\cos x &{}+48\overset{\prime\prime}{,}391\sin x\\
&+14{,}295\cos2x &{}-36{,}951\sin2x\\
&+16{,}161\cos3x &{}-8{,}251\sin3x\\
&+6{,}273\cos4x &{}+2{,}923\sin4x\\
&+0{,}162\cos5x &{}+6{,}354\sin5x\\
&-0{,}674\cos6x &{}+3{,}129\sin6x\\
&+1{,}702\cos7x &{}+0{,}899\sin7x\\
&+0{,}922\cos8x &{}-0{,}065\sin8x\\
&+0{,}460\cos9x &{}-0{,}414\sin9x\\
&+0{,}181\cos10x &{}-0{,}326\sin10x\\
&+0{,}063\cos11x &{}-0{,}198\sin11x\\
&+0{,}023\cos12x &{}-0{,}101\sin12x
\end{aligned}
\]
\[
\begin{aligned}
\text{2) for }Q={}&-544\overset{\prime\prime}{,}924\\
&-98{,}084\cos x &{}-6\overset{\prime\prime}{,}178\sin x\\
&+27{,}589\cos2x &{}-10.928\sin2x\\
&+8{,}342\cos3x &{}+4{,}711\sin3x\\
&-0{,}875\cos4x &{}+1{,}294\sin4x\\
&-4{,}480\cos5x &{}-1{,}766\sin5x\\
&-1{,}813\cos6x &{}-1{,}284\sin6x\\
&-0{,}179\cos7x &{}+1{,}643\sin7x\\
&+0{,}434\cos8x &{}+1{,}020\sin8x\\
&+0{,}414\cos9x &{}+0{,}460\sin9x\\
&+0{,}326\cos10x &{}+0{,}181\sin10x\\
&+0{,}198\cos11x &{}+0{,}063\sin11x\\
&+0{,}101\cos12x &{}+0{,}023\sin12x
\end{aligned}
\]
\noindent\hspace*{\parindent}3) for \(\log N\) and \(\psi\), so that
\[
P=N\cos\psi,\qquad Q=N\sin\psi.
\]
\noindent\(\psi\) is ultimately required to \textls[100]{arcseconds}, and \(\log N\) to 5 decimal places. Ideally,
one would have this table for \(x=0\), \(x=1^0\), \(x=2^0\), \(x=3^0\) etc. up to
\(x=360^0\), but that would be far too formidable a task. We
\SourceLeaf{573}
must therefore set narrower limits. I think we may \textls[100]{if need be} be content
to let the values of \(x\) increase from \(15^0\) to \(15^0\), so that
24 such calculations have to be made. That is how I would like the trial calculation done.
Several shortcuts for particular cases can be used here,
which you will readily find for yourself and which the following identities
\[
\left.\begin{aligned}
\sin15^0&=\BareRoot{\tfrac{3}{8}}-\BareRoot{\tfrac{1}{8}}\\
\sin45^0&=2\,\mathord{.}\,\BareRoot{\tfrac{1}{8}}
&\qquad\sin60^0&=\BareRoot{\tfrac{3}{4}}\\
\sin75^0&=\BareRoot{\tfrac{3}{8}}+\BareRoot{\tfrac{1}{8}}
\end{aligned}\right\}\ \text{enable.}
\]
\noindent\hspace*{\parindent}.\ .\ .\ .\ .
\clearpage
\endgroup
